\subsection{映射的典范扩张}

设 \(S = (c_1, c_2, \ldots, c_m)\) 是一个具有 \(n\) 个项的阶梯构造方案。设 \(E_1, \ldots, E_n, E_1', \ldots, E_n'\) 为集合（即 \(\mathcal{G}\) 中的项），并设 \(f_1, \ldots, f_n\) 为 \(\mathcal{G}\) 中的项，使关系“ \(f_i\) 是从 \(E_i\) 到 \(E_i'\) 的映射”在 \(\mathcal{G}\) 中对每个 \(1 \leqslant i \leqslant n\) 都是定理。设 \(A_1, \ldots, A_m\) （分别地， \(A_1', \ldots, A_m'\) ) 为方案 S 在 \(E_1, \ldots, E_n\) （分别地， \(E_1', \ldots, E_n'\) ) 上的阶梯构造。我们按下述条件逐步定义一个由 \(m\) 个项 \(g_1, \ldots, g_m\) 组成的序列，使 \(g_i\) 是从 \(A_i\) 到 \(A_i'\) 的映射（ \(1 \leqslant i \leqslant m\) ）：

\begin{enumerate}
\def\labelenumi{(\alph{enumi})}
\item
  如果 \(c_{i} = (0, b_{i})\) ，因此 \(A_{i} = E_{b_{i}}\) 并且 \(A_{i}' = E_{b_{i}}'\) ，那么 \(g_{i}\) 是映射 \(f_{b_{i}}\) .
\item
  如果 \(c_{i} = (a_{i}, 0)\) ，因此 \(\mathbf{A}_{i} = \mathfrak{P}(\mathbf{A}_{a_{i}})\) 并且 \(\mathbf{A}_{i}' = \mathfrak{P}(\mathbf{A}_{a_{i}}')\) ，那么 \(g_{i}\) 是典范扩张 \(\hat{g}_{a_i}\) ，即 \(g_{a_i}\) 到子集族的扩张（第二章，第5节，第1号）。
\item
  如果 \(c_{i} = (a_{i}, b_{i})\) ，其中 \(a_{i} \neq 0\) 并且 \(b_{i} \neq 0\) ，因此
\end{enumerate}

\[
\mathrm{A} _ {i} = \mathrm{A} _ {a _ {i}} \times \mathrm{A} _ {b _ {i}} \quad \text { 且 } \quad \mathrm{A} _ {i} ^ {\prime} = \mathrm{A} _ {a _ {i}} ^ {\prime} \times \mathrm{A} _ {b _ {i}} ^ {\prime},
\]

则 \(g_{i}\) 是典范扩张 \(g_{a_i} \times g_{b_i}\) ，它将 \(g_{a_i}\) 和 \(g_{b_i}\) 扩张到 \(A_{a_i} \times A_{b_i}\) 上（第二章，第3节，第9号）。

该序列的最后一项 \(g_{m}\) 称为映射 \(f_{1}, \ldots, f_{n}\) 的具有方案S的典范扩张，并将用 \(\langle f_{1}, \ldots, f_{n} \rangle^{S}\) 表示。以下准则可以逐步验证：

CST1. 如果 \(f_{i}\) 是从 \(\mathbf{E}_i\) 到 \(\mathbf{E}_i'\) 的映射，并且如果 \(f_{i}'\) 是从 \(\mathbf{E}_i'\) 到 \(\mathbf{E}_i''\) 的映射 ( \(1 \leqslant i \leqslant n\) )，那么对于每个 \(n\) 项上的层级构造方案，我们有

\[
\langle f _ {1} ^ {\prime} \circ f _ {1}, f _ {2} ^ {\prime} \circ f _ {2}, \dots , f _ {n} ^ {\prime} \circ f _ {n} \rangle^ {\mathrm{S}} = \langle f _ {1} ^ {\prime}, f _ {2} ^ {\prime}, \dots , f _ {n} ^ {\prime} \rangle^ {\mathrm{S}} \circ \langle f _ {1}, f _ {2}, \dots , f _ {n} \rangle^ {\mathrm{S}}.
\]

CST2. 若 \(f_{i}\) 对 \(1 \leqslant i \leqslant n\) 是单射（相应地，满射），那么 \(\langle f_{1}, \ldots, f_{n} \rangle^{S}\) 是单射（相应地，满射）。

该判据由扩展的相应性质得出。 \(\hat{g}\) (第二章，§5，第1号，命题1)及扩展 \(g \times h\) （第二章，第3节，第9号）。

CST3. 如果 \(f_{i}\) 是 \(\mathbf{E}_i\) 到 \(\mathbf{E}_i'\) 上的双射，并且如果 \(f_{i}^{-1}\) 是(*)的逆双射，则 \(\langle f_1, \ldots, f_n \rangle^{S}\) 是一个双射且 \(\langle f_1^{-1}, \ldots, f_n^{-1} \rangle^{S}\) 是其逆；换言之，

\[
(\langle f _ {1}, \dots , f _ {n} \rangle^ {\mathrm{S}}) ^ {- 1} = \langle f _ {1} ^ {- 1}, \dots , f _ {n} ^ {- 1} \rangle^ {\mathrm{S}}.
\]

这直接从CST1和CST2得出。

